\section{Problem Formulation}\label{chap:problem}
\subsection{Existence and Uniqueness for a constrained optimization problem}
To prove the existence and uniqueness of a minimizer for a constrained optimization problem,
three properties are required. Feasibility, Coercivity, and Convexity. Feasibility ensures that the constraint set is non-empty, meaning
that there is at least one point that satisfies all constraints. Coercivity of the objective
function guarantees that the function tends to infinity as the norm of the variable tends to infinity,
ensuring the existence of a global minimum on the feasible set. Finally, convexity of the objective
function implies that any local minimizer is also a global one, and if the function is strictly convex,
it ensures that the minimizer is unique. Therefore, the combination of feasibility, coercivity, and
strict convexity, establishes both the existence and uniqueness of a global minimizer.

\subsection{The Primal Problem (P)}
We begin with the first soft-margin classification problem we are interested in solving, which is
formulated as a convex optimization problem. This will serve as the primal problem in our analysis.

\begin{align*}\label{eq:primal-problem}
	\min_{\symbf{w}, b, \symbf{\xi}} \quad
	 & f(\symbf{w}, b, \symbf{\xi})
	= \frac{1}{2}\|\symbf{w}\|_2^2 + C\,\symbf{1}_M^\top \symbf{\xi}
	\quad
	\text{subject to}
	\quad
	\begin{cases}
		\symbf{1}_M - Y\bigl(X\,\symbf{w} + b\,\symbf{1}_M\bigr) - \symbf{\xi} & \le \symbf{0}_M, \\[6pt]
		\symbf{\xi} \;\;                                                       & \ge \symbf{0}_M,
	\end{cases}\tag{P}
\end{align*}

Here, $\symbf{w} \in \R^d$ represents the weight vector, $b \in \R$ is the bias term, and $\symbf{\xi} \in \R^M$ is the vector of slack variables.
The data matrix is $X = [\symbf{x}_1, \ldots, \symbf{x}_M] \in \R^{d \times M}$ with labels $\symbf{y}=[y_1, \ldots, y_M]^\top \in \{\pm1\}^M$ and diagonal label matrix $Y = \operatorname{diag}(y_1,\ldots,y_M)$.
The parameter $C>0$ controls the penalty for misclassification.

\begin{theorem}{Existence \& Uniqueness of Primal Solution}{existence-uniqueness-primal}
	The primal problem \eqref{eq:primal-problem} always admits at least one global minimizer. Moreover, the weight vector $\symbf{w}^\star$ is unique due to strict convexity of the objective in $\symbf{w}$. However, the bias $b^\star$ and slack variables $\symbf{\xi}^\star$ may fail to be unique under certain degenerate configurations.
\end{theorem}

\begin{proof}[Proof of Theorem~\ref{thm:existence-uniqueness-primal}]
	\paragraph{Existence of a solution.}
	The primal problem~\eqref{eq:primal-problem} is a convex quadratic program with affine inequality constraints. To establish existence of a minimizer, we verify feasibility and coercivity of the objective.

	\begin{itemize}
		\item \textbf{Feasibility:} Let $\symbf{w} = \symbf{0}_d$, $b = 0$, and $\symbf{\xi} = \symbf{1}_M$. Then,
		      \[
			      \symbf{1}_M - Y\left(X\symbf{w} + b\,\symbf{1}_M\right) - \symbf{\xi}
			      = \symbf{1}_M - Y(\symbf{0}_M) - \symbf{1}_M = \symbf{0}_M \le \symbf{0}_M,
		      \]
		      and clearly $\symbf{\xi} = \symbf{1}_M \ge \symbf{0}_M$. Thus, this choice satisfies all constraints, and the feasible region is non-empty.

		\item \textbf{Coercivity:} The objective function
		      \[
			      f(\symbf{w}, b, \symbf{\xi}) = \frac{1}{2} \|\symbf{w}\|_2^2 + C\, \symbf{1}_M^\top \symbf{\xi}
		      \]
		      is coercive in $\symbf{w}$ and $\symbf{\xi}$. As either $\|\symbf{w}\| \to \infty$ or $\|\symbf{\xi}\| \to \infty$, the objective tends to infinity. Since the feasible region is non-empty and convex, and the objective is continuous and coercive, the problem attains a minimum.
	\end{itemize}

	\paragraph{Uniqueness of a solution.}
	\begin{itemize}
		\item \textbf{Uniqueness of $\symbf{w}^\star$:}
		      The term $\tfrac{1}{2}\|\symbf{w}\|^2$ is strictly convex in $\symbf{w}$, while the remaining terms are affine. As a result, the full objective is strictly convex in $\symbf{w}$, guaranteeing that $\symbf{w}^\star$ is unique.

		\item \textbf{Possible non-uniqueness of $b^\star$:}
		      The variable $b$ appears only linearly in the constraints and not in the objective. In cases where the constraints do not uniquely determine $b$—for example, when multiple values of $b$ yield the same classification margin—a degree of freedom in $b$ remains. This can happen, for instance, when there are no support vectors exactly on the margin.

		\item \textbf{Possible non-uniqueness of $\symbf{\xi}^\star$:}
		      The slack variables $\xi_i$ only appear linearly in both the objective and the constraints. For points that are not support vectors (i.e., not active in the constraint), the optimal $\xi_i^\star$ may be any non-negative value satisfying the constraint, leading to potential non-uniqueness. However, for support vectors that lie within or on the margin, $\xi_i^\star$ becomes uniquely determined.
	\end{itemize}
\end{proof}

\subsection{The Dual Problem (DP)}

The second soft-margin classification problem we are interested in solving, which will be refered
as the dual problem, is also a convex optimization problem:

\begin{align}\label{eq:dual-problem}
	\min_{\symbf{\alpha}}
	\;\;\frac{1}{2}\,\inner*{\symbf{\alpha},\,Y\,G\,Y\,\symbf{\alpha}}
	\;-\;\inner*{\symbf{1}_M,\;\symbf{\alpha}}
	\quad
	\text{subject to}
	\quad
	\begin{cases}
		\inner{\symbf{y}}{\symbf{\alpha}} = 0, \\
		\symbf{0}_M \;\;\leq\;\; \symbf{\alpha} \;\;\leq\;\; C\,\symbf{1}_M.
	\end{cases}
\end{align}
Here $G = \bigl(\inner*{\symbf{x}_i}{\symbf{x}_j}\bigr)_{i,j=1}^M$ is the Gram matrix, and $Y=\operatorname{diag}(y_1,\dots,y_M)$.

\begin{theorem}{Existence \& Uniqueness of Dual Solution}{existence-uniqueness-dual}
	The dual problem \eqref{eq:dual-problem} always admits at least one global minimizer. Moreover, if $YGY$ is strictly positive-definite on the subspace $\{\symbf{\alpha}\mid \inner{\symbf{y}}{\symbf{\alpha}}=0\}$, then the solution $\symbf{\alpha}^\star$ is unique. In degenerate cases, multiple optimal solutions may exist.
\end{theorem}

\begin{proof}
	\paragraph{Existence of a solution.}
	The dual problem \eqref{eq:dual-problem} is a convex quadratic program in $\symbf{\alpha}$, subject to a linear equality constraint $\langle \symbf{y}, \symbf{\alpha} \rangle = 0$ and box constraints $0 \leq \alpha_i \leq C$ for all $i = 1, \ldots, M$. We define the feasible region as:
	\[
		\Omega = \left\{ \alpha \in \mathbb{R}^M \;\middle|\; 0 \leq \alpha_i \leq C, \;\langle \symbf{y}, \alpha \rangle = 0 \right\},
	\]
	which is the intersection of the closed box $[0, C]^M$ with a hyperplane. This intersection is nonempty: because the dataset includes both positive and negative labels (i.e., $\exists\,p,n$ such that $y_p = 1$, $y_n = -1$), one can construct a feasible point by choosing $\alpha_p = \alpha_n$ and setting all other components to zero. Hence, $\Omega$ is closed, bounded, and nonempty.

	The dual objective function,
	\[
		f(\alpha) = \frac{1}{2} \langle \alpha, YGY \alpha \rangle - \langle \mathbf{1}_M, \alpha \rangle,
	\]
	is continuous and convex, being a quadratic form with symmetric matrix $YGY$ and linear term $- \mathbf{1}_M^\top \alpha$. Since $f$ is continuous and convex over a compact feasible set $\Omega$, the extreme value theorem guarantees the existence of a global minimizer $\alpha^\star$.

	\paragraph{Uniqueness of a solution.}
	The Hessian of $f$ is $\nabla^2 f(\alpha) = YGY$, which is symmetric and positive semi-definite because $G$ is a Gram matrix and $Y$ is a diagonal matrix with $\pm 1$ entries. If $YGY$ is strictly positive-definite on the affine subspace $\{\alpha \in \mathbb{R}^M : \langle y, \alpha \rangle = 0\}$, then $f$ is strictly convex on the feasible region $\Omega$, and the minimizer $\alpha^\star$ is unique. In contrast, if $YGY$ is only positive semi-definite (e.g., if $G$ is rank-deficient), then flat directions may exist in the objective, and multiple global minimizers may be possible.
\end{proof}


% \section{Existence and Uniqueness of the Soft-Margin SVM Problem}
% \[
% \min_{w \in \mathbb{R}^d,\;b \in \mathbb{R},\;\xi \in \mathbb{R}^M}
% \quad \frac{1}{2}\,\|w\|^2 \;+\; C \sum_{i=1}^M \xi_i
% \quad\text{subject to}\quad
% \begin{cases}
% y_i\bigl(\langle w, x_i\rangle + b\bigr)\;\ge\;1 \;-\;\xi_i,\\[4pt]
% \xi_i \;\ge\; 0,
% \end{cases}
% \quad i = 1,\dots,M,
% \]

% \[
% \min_{\alpha \in \mathbb{R}^M} 
% \frac{1}{2} \langle \alpha, YGY \alpha \rangle 
% - \langle \mathbf{1}_M, \alpha \rangle
% \quad \text{subject to} \quad
% \begin{cases}
% \langle y, \alpha \rangle = 0,\\[6pt]
% 0 \le \alpha_i \le C
% \end{cases}
% \quad \text{for } i = 1, \dots, M.
% \]

% where each $y_i\in\{-1,+1\}$ and $C>0$ is a fixed penalty parameter.

% To prove the existence and uniqueness of a minimizer for a constrained optimization problem, 
% three properties are required. Feasibility ensures that the constraint set is non-empty, meaning 
% that there is at least one point that satisfies all constraints. Coercivity of the objective 
% function guarantees that the function tends to infinity as the norm of the variable tends to infinity,
% ensuring the existence of a global minimum on the feasible set. Finally, convexity of the objective 
% function implies that any local minimizer is also a global one, and if the function is strictly convex,
% it ensures that the minimizer is unique. Therefore, the combination of feasibility, coercivity, and 
% strict convexity, establishes both the existence and uniqueness of a global minimizer.


% \subsection*{1.\ Non-emptiness and Feasibility}

% The feasible set of $(w,b,\{\xi_i\})$ is non-empty. For example, one can always pick
% \[
% w = 0,\quad b = 0,\quad \xi_i = 1 \quad \forall\,i,
% \]
% so that
% \(
% y_i \langle w,x_i\rangle + b = 0
% \)
% and
% \(\,1 - \xi_i \le 0,\)
% thus satisfying
% $y_i(\langle w,x_i\rangle + b)\ge 1-\xi_i.$
% Hence there is at least one feasible point, guaranteeing the set of all feasible triples $(w,b,\xi)$ is non-empty.

% \subsection*{2.\ Coercivity and Existence of a Global Minimizer}

% Observe that the objective
% \[
% \frac{1}{2}\|w\|^2 \;+\; C\sum_{i=1}^M \xi_i
% \]
% is \emph{coercive} in the variables $(w,\xi)$: as $\|w\|\to\infty$ or $\sum_i \xi_i \to \infty,$ the cost grows unbounded. Together with the linear constraints $\xi_i\ge 0$ and the fact that the feasible set is not empty, standard results in convex analysis imply the existence of \emph{some} global minimizer $(w^*,b^*,\xi^*)$.

% \subsection*{3.\ Convexity and (Possible) Uniqueness}

% The function 
% \[
% \frac12 \|w\|^2 + C \sum_{i=1}^M \xi_i
% \]
% is convex in $(w,\xi)$, and \emph{affine} (hence convex) in $b$. Therefore, the entire objective is convex on a convex feasible region (defined by affine/linear inequalities). This guarantees that every local minimizer is also a global minimizer.

% \paragraph{Uniqueness.} 
% Strict convexity in the $w$-part implies that, \emph{for fixed $b$ and $\{\xi_i\}$}, the objective in $w$ is strictly convex, which typically ensures a unique $w^*$ \emph{conditional} on $b,\{\xi_i\}.$ However, because the objective is \emph{linear} in $b$, strict convexity does not extend to the $(b,w,\xi)$ space as a whole. In many cases, $b^*$ can still be pinned down by the margin constraints, but in some situations there may be multiple solutions that share the same minimal cost. 

% Hence, if the data are in ``general position,'' one typically obtains a \emph{unique} $w^*$ and a well-defined $b^*$, but strictly speaking, the problem is not globally \emph{strictly} convex in \emph{all} variables $(w,b,\xi)$ simultaneously. Therefore, uniqueness may or may not hold in full, depending on degeneracies in the data or constraints.


% \[
% f(\alpha) := \min_{\alpha \in \mathbb{R}^M} \frac{1}{2} \langle \alpha, YGY \alpha \rangle - \langle \mathbf{1}_M, \alpha \rangle \quad \text{s.t.} \quad 
% \begin{cases}
% \langle y, \alpha \rangle = 0 \\
% 0 \leq \alpha_i \leq C \quad \text{for } i = 1, \ldots, M
% \end{cases}
% \]

% \[
% Y = \text{diag}(y_1, \ldots, y_M) \quad \text{where } y_i = 1 \vee -1 \text{ for } i = 1, \ldots, m
% \]

% \[
% G = (\langle x_i, x_j \rangle)_{i,j=1,\ldots, M}, \quad \text{thus } G \text{ is the kernel matrix and SPD}
% \]

% \[
% \mathbf{1} = (1, \ldots, 1)^T
% \]

% \subsection*{1. Feasibility:}
% The area where the function $f(\alpha)$ can exist is given by
% \[
% \Omega = \left\{ \alpha \in \mathbb{R}^M \mid 0 \leq \alpha_i \leq C, \, \langle y, \alpha \rangle = 0 \right\}
% \]
% Thus the feasible region is the intersection of the area $[0, C]^M$ with the area $\sum_i y_i \alpha_i = 0$.

% Because there must at least be one positive and negative label, two indices, $p$ and $n$, can be picked such that $y_p = 1$ and $y_n = -1$ while $\alpha_p = \alpha_n$, letting all other $\alpha$ values be zero, satisfying both constraints. Guaranteeing that the feasible set is not zero.

% Also, $\alpha_n$ and $\alpha_p$ can be set to $0$ or $C$, thus boundary and corner points are also feasible.

% Thus $\Omega$ is a non-empty, closed and bounded area.

% \subsection*{2. Convexity:}
% Rewriting $f(\alpha) = \frac{1}{2} \langle \alpha, YGY \alpha \rangle - \langle \mathbf{1}_M, \alpha \rangle$ to

% \[
% f(\alpha) = \frac{1}{2} \alpha^T YGY \alpha - \mathbf{1}^T \alpha
% \]

% Finding the Hessian of this we find:

% \[
% \nabla f(\alpha) = Y G Y \alpha - \mathbf{1} \quad \Rightarrow \quad \nabla^2 f(\alpha) = Y G Y
% \]

% Where $G$ is semi-positive definite by definition as it is a gram matrix, and $Y$ is a diagonal matrix with $\pm 1$ values. Thus $YGY$ is symmetric and PSD, meaning $f$ is convex.
